\documentclass[11pt]{article}
\usepackage[margin=1in]{geometry}
\usepackage{amsmath,amssymb,amsthm,mathtools}
\usepackage[colorlinks=true,linkcolor=blue,citecolor=blue,urlcolor=blue]{hyperref}
\newtheorem{theorem}{Theorem}
\newtheorem{proposition}[theorem]{Proposition}
\newcommand{\C}{\mathbb C}
\newcommand{\R}{\mathbb R}
\newcommand{\D}{\mathcal D}
\newcommand{\Ba}{\mathcal B^a}
\DeclareMathOperator{\diag}{diag}
\DeclareMathOperator{\ev}{ev}
\title{Forward generation does not ensure GNS descent:\\an explicit continuous counterexample}
\author{Alec Kriebel\\Independent researcher\\
\small\href{https://orcid.org/0009-0001-9320-500X}{ORCID: 0009-0001-9320-500X}}
\date{13 October 2026}
\begin{document}
\maketitle
\begin{abstract}
For a general conditional-expectation dilation of a Markov semigroup,
generation by the observed algebra and its forward translates does not
ensure descent through the expectation's fixed GNS representation. We give
a unital separable, injective, point-norm continuous example for the
two-state exponential relaxation semigroup. Its GNS image is $M_2(\C)$,
and at every positive time an explicit nonzero positive element leaves
the representation kernel. This refutes the concrete generation-only
proposal in Skeide's 2009 module-approach problem. The clock/cocycle
method is established: the negative phenomenon is also a short corollary
of an earlier clock dilation of Burgarth, Facchi and Hillier. The note
makes the connection explicit and supplies a small matrix realization
with exact witnesses; it makes no claim of literature-first priority.
\end{abstract}

\paragraph{Disclosure.}
AI tools were used extensively in constructing, checking, auditing and
writing this work. This is an unrefereed research note; its adversarial
reviews were performed by AI tools and do not constitute human peer review.

\section*{The question and the result}
In the first module-approach problem of \cite[pp.~542--543]{owr}, Skeide
proposes passing a general dilation $(A,\theta,i,p)$ through the GNS
correspondence of $p$, and then extending its represented dynamics to all
adjointable operators. Generation of $A$ by $i(B)$ and its time translates
is offered as a sufficient ``not too big'' condition. We use the report's
general definition \cite[pp.~497--498]{owr}: $ip$ is a conditional
expectation onto $i(B)$ and $p\theta_t i=T_t$. The printed direction of
$p$ in the problem is reversed; the definition and diagrams specify
$p:A\to B$. The stronger equation $p\theta_t=T_t p$ on all of $A$ is
not assumed.

Let $\D\subset M_2(\C)$ be the diagonal algebra, identified with $\C^2$,
and let $\Delta:M_2(\C)\to\D$ extract the diagonal.
\begin{theorem}\label{thm:main}
There exist a unital separable C*-algebra $A$, a point-norm continuous
semigroup $(\theta_t)_{t\geq0}$ of injective unital *-endomorphisms,
a unital embedding $i:\D\to A$, and a conditional expectation
$ip:A\to i(\D)$ such that
\[
 A=C^*(\theta_r(i(\D)):r\geq0),\qquad
 T_t=p\theta_t i,\qquad
 P_t=\frac12\begin{pmatrix}1+e^{-t}&1-e^{-t}\\1-e^{-t}&1+e^{-t}\end{pmatrix}
\]
is the transition matrix of $T_t$ on $\C^2$.
For the left representation $\pi$ in the GNS correspondence of $p$,
\[
 \pi(A)\cong M_2(\C),\qquad
 \theta_s(\ker\pi)\not\subseteq\ker\pi\quad\text{for every }s>0.
\]
Each failure has an explicit nonzero positive witness $a_s\in\ker\pi$.
Consequently no semigroup on that represented algebra, or on all
$\Ba(E)$ of the fixed GNS correspondence, can intertwine $\pi$ and $\theta$.
\end{theorem}

\section*{Construction and proof}
Let $X=[-\infty,+\infty]$ be the two-endpoint compactification of $\R$
with its order topology. Define
\[
 g(x)=\begin{cases}0,&x\leq0,\\\tfrac12\arccos(e^{-x}),&x>0,\end{cases}
 \qquad g(-\infty)=0,\quad g(+\infty)=\pi/4,
\]
and put
\[
 R(z)=\begin{pmatrix}\cos z&-\sin z\\\sin z&\cos z\end{pmatrix},
 \qquad V(x)=R(g(x)),\qquad u_t(x)=V(x)^*V(x+t).
\]
Translation fixes the endpoints of $X$. It is jointly continuous in
$(t,x)$: for time in a bounded neighborhood, sufficiently far tails
remain in any prescribed endpoint neighborhood. Thus, on $C(X,M_2(\C))$,
\[
 \alpha_t(f)(x)=u_t(x)f(x+t)u_t(x)^*,\qquad t\in\R,
\]
defines a point-norm continuous automorphism group. Indeed
$u_s(x)u_t(x+s)=u_{s+t}(x)$, and joint continuity on compact time
intervals times $X$ gives uniform continuity for each fixed $f$.
No differentiability of $g$ at zero is required.

View $\D$ as the constant diagonal functions and set
\[
 A=C^*(\alpha_r(\D):r\geq0),\quad
 \theta_t=\alpha_t|_A\ (t\geq0),\quad i(b)(x)=b,\quad
 p(f)=\Delta(f(0)).
\]
The shift of a generator at $r$ is the generator at $r+t$, so $A$ is
forward invariant and the asserted generation equality is exact.
Restrictions of the ambient automorphisms are injective unital
endomorphisms. Norm continuity allows rational times to generate $A$,
proving separability. Evaluation and $\Delta$ are unital completely
positive; $pi=\mathrm{id}_{\D}$ and
$p(i(b_1)fi(b_2))=b_1p(f)b_2$. Hence $ip$ is a conditional expectation.

At $x=0$, $u_t(0)=R(g(t))$ for $t\geq0$. Compressing a constant
diagonal matrix gives weights
$\cos^2 g(t)=(1+e^{-t})/2$ and $\sin^2 g(t)=(1-e^{-t})/2$.
These are precisely $P_t$. Writing
\[
 J=\tfrac12\begin{pmatrix}1&1\\1&1\end{pmatrix},\qquad
 K=\tfrac12\begin{pmatrix}1&-1\\-1&1\end{pmatrix},\qquad P_t=J+e^{-t}K,
\]
we have $J^2=J$, $K^2=K$, $JK=KJ=0$ and $J+K=I$.
This proves the Markov semigroup law for every real $t\geq0$.
The maps are UCP and uniformly continuous on $\C^2$.

\paragraph{The actual GNS representation.}
Let $q=\diag(1,0)$ and $s_0=\log2$. Evaluation of
$\theta_{s_0}(q)$ at zero is
\[
 Q=\begin{pmatrix}3/4&\sqrt3/4\\\sqrt3/4&1/4\end{pmatrix}.
\]
The evaluation image contains $q$, $1-q$, $qQ(1-q)=(\sqrt3/4)e_{12}$
and its adjoint. Therefore $\ev_0(A)=M_2(\C)$.
The GNS correspondence is concretely
\[
 E=M_2(\C),\qquad z\cdot b=zb,\qquad
 \langle z,w\rangle=\Delta(z^*w),\qquad \xi=I_2,\qquad
 \pi(f)z=f(0)z.
\]
To verify this, the map $a\otimes b\mapsto a(0)b$ on the algebraic
GNS construction preserves its inner product. The concrete inner
product is definite because its diagonal entries are column squared
norms, and the evaluation image proves surjectivity and cyclicity.
Finite dimension gives completeness. In particular,
\[
 \ker\pi=\{f\in A:f(0)=0\}.
\]
This is the full representation kernel, not an identification with a
single vector-state null space in general. The two columns identify
$\Ba(E)=M_2(\C)\oplus M_2(\C)$; $\pi(A)$ is the diagonal copy.
The module is full since $\langle\xi,\xi\rangle=1_{\D}$.

\paragraph{A positive witness at every positive time.}
Fix $s>0$, put $r=e^{-s}\in(0,1)$ and $Q_s=\theta_s(q)(0)$, and define
\[
 a_s=q\theta_s(q)q-\frac{1+r}{2}q.
\]
Then $a_s\in A$ and $a_s(0)=0$. At $x=s$ its coefficient of $q$ is
\begin{align*}
 d_s&=\cos^2(g(2s)-g(s))-\frac{1+r}{2}\\
 &=\frac{r^3+(1-r^2)\sqrt{1+r^2}-r}{2}
 =\frac{(1-r^2)(\sqrt{1+r^2}-r)}{2}>0.
\end{align*}
Here the angle-difference identity uses
$\cos(2g(s))=r$, $\cos(2g(2s))=r^2$, and their positive sines.
Consequently
\[
 \theta_s(a_s)(0)=d_s Q_s\ne0.
\]
Moreover $a_s$ is positive everywhere. For $x>0$,
$g'(x)=1/(2\sqrt{e^{2x}-1})$ is positive and decreasing. Comparing
its integrals on $[x,x+s]$ and $[0,s]$ gives
$0\leq g(x+s)-g(x)\leq g(s)$. For $x\leq0$ the same bound follows
from monotonicity and $g(x)=0$, and at the endpoints the increment is
zero. Since $\cos^2$ decreases on $[0,\pi/4]$,
\[
 a_s(x)=\left(\cos^2(g(x+s)-g(x))-\cos^2g(s)\right)q\geq0.
\]
Strict positivity of $d_s$ proves $a_s\ne0$. At $s=\log2$,
$d_s=3(\sqrt5-1)/16$, giving an exact radical witness.
An intertwining map would send both $\pi(0)$ and $\pi(a_s)$ to the
same value while their required images are $0$ and $d_s Q_s$.
This proves the theorem.

\section*{Antecedents and the finite-image obstruction}
The point-localized clock and translated-unitary construction are
established in the cited arXiv v2 of Burgarth--Facchi--Hillier (BFH)
\cite[Definition~4.1, Theorem~4.4, Step~2, equation~(11)]{bfh}.
Finite time-dependent unitary dilations and their initial square-root
singularities are treated by Dive--Mintert--Burgarth
\cite[Method, Appendices~A--B, arXiv~v3]{dmb}. Their first-derivative
obstruction has a bounded-total-Hamiltonian formulation in the cited
arXiv v3 of \cite[Theorem~1]{ende}.
The following elementary conditional-expectation version explains the
connection; it is proved here rather than attributed to those papers
as a fixed-GNS theorem.

\begin{proposition}\label{prop:finite}
For a point-norm continuous general unital dilation of a norm-continuous
Markov semigroup on a finite-dimensional algebra $B$, if $\pi(A)$ is
finite dimensional and $\ker\pi$ is invariant at every time, then the
generator of the compressed semigroup is a *-derivation on $B$.
\end{proposition}
\begin{proof}
Kernel invariance gives a continuous endomorphism semigroup $\beta_t$
on $C=\pi(A)$. Finite dimension makes continuity operator-norm
continuity. For small $h>0$, $S=\int_0^h\beta_u\,du$ is invertible
because $S/h\to I$. The identity
$\beta_tS=\int_t^{t+h}\beta_u\,du$ proves right differentiability
at zero. Its generator $\delta$ is a *-derivation by differentiation
of multiplicativity. The embedding $j=\pi i$ is faithful, and
$\bar p(c)=\langle\xi,c\xi\rangle$ is a $B$-bimodule UCP map
with $\bar p j=\mathrm{id}_B$. Thus, identifying $B$ with $j(B)$,
\[
 L(bc)=\bar p(\delta(bc))=\bar p(\delta(b))c+b\bar p(\delta(c))
       =L(b)c+bL(c).
\]
Adjoints are preserved as well.
\end{proof}

In fact BFH's earlier explicit amplitude-damping clock already yields
the negative generation-only conclusion as a short corollary. The model
is in Section~2, equations~(1)--(3), of the cited arXiv v2 of \cite{bfh}.
Its finite unitary path has angle
\[
 G(x)=\operatorname{sgn}(x)\arccos(e^{-|x|/2}),\qquad G(0)=0,
\]
with finite endpoint limits.
It therefore extends to $[-\infty,+\infty]$, where its translated
cocycle is point-norm continuous. Restriction to the forward-generated
observable algebra preserves the conditional expectation and the
compression. Evaluation followed by finite bath compression gives a
finite-dimensional GNS image. The decaying-level projection satisfies
$T_t(q)=e^{-t}q$, whereas a derivation satisfies $qL(q)q=0$.
Proposition~\ref{prop:finite} rules out invariance at all times.
The generic strong continuity of BFH's Hilbert-space unitaries alone
would not justify point-norm continuity on all $C_b(\R)$; the endpoint
check for this explicit path is essential.

Thus the general negative phenomenon is not a new clock-dilation
mechanism. Theorem~\ref{thm:main} gives a small $M_2/\C^2$ realization
and positive witnesses at every positive time. Bounded primary-source
searches through 13 October 2026 found no publication explicitly stating
the exact OWR fixed-GNS counterexample, but neither that search nor the
explicit refinement certifies first priority. Full source locations,
hashes and the prior-result deduction are in the verification package.

\section*{Scope and reproducibility}
This answers negatively the concrete forward-generation-only proposal
in \cite{owr}. Failure of its first step refutes the proposed two-step
conjunction. It also defeats any sufficient condition logically implied
by forward generation alone. Useful stronger or incomparable minimality
conditions, and extension after kernel invariance is separately imposed,
remain outside the result. The related fixed-action question in
\cite[Section~5]{ss} should be read with these same limits.

The example is a general C*-dilation with a nonfaithful expectation.
It is not strong: $p(a_s)=0$ but $p\theta_s(a_s)\ne0$. It is not a
weak/corner dilation, and no normal von Neumann version or nonexistence
of a different dilation of the same Markov semigroup is asserted.
These categories are distinguished in the cited arXiv v3 of
\cite{shalit}: Definition~2.3, diagram~(2.3), and Appendix~R(i),
pp.~211--212. The locator refers specifically to that preprint version.

The accompanying package contains this source and PDF, the frozen
candidate, independent exact arithmetic checks, analytic audits,
priority evidence and a hash-checked replay program. Finite computations
supplement the all-real-time proof; they do not establish its continuous
quantifiers. Hashes certify byte integrity, not mathematical truth.
No third-party full texts are redistributed.

\begin{thebibliography}{9}
\bibitem{owr} M. Skeide, Some Problems, in \emph{Mini-Workshop:
Product Systems and Independence in Quantum Dynamics}, Oberwolfach
Reports \textbf{6} (2009), 493--548, especially 542--543.
\href{https://doi.org/10.4171/OWR/2009/09}{doi:10.4171/OWR/2009/09}.
\bibitem{bfh} D. Burgarth, P. Facchi and R. Hillier,
Control of quantum noise: on the role of dilations,
Ann. Henri Poincar\'e \textbf{24} (2023), 325--347.
\href{https://doi.org/10.1007/s00023-022-01211-y}{doi:10.1007/s00023-022-01211-y};
\href{https://arxiv.org/abs/2010.00279v2}{arXiv:2010.00279v2}.
\bibitem{dmb} B. Dive, F. Mintert and D. Burgarth,
Quantum simulations of dissipative dynamics: Time dependence instead of size,
Phys. Rev. A \textbf{92} (2015), 032111.
\href{https://doi.org/10.1103/PhysRevA.92.032111}{doi:10.1103/PhysRevA.92.032111};
\href{https://arxiv.org/abs/1503.01010v3}{arXiv:1503.01010v3}.
\bibitem{ende} F. vom Ende, Quantum-Dynamical Semigroups and the
Church of the Larger Hilbert Space, Open Syst. Inf. Dyn. \textbf{30}
(2023), 2350003.
\href{https://doi.org/10.1142/S1230161223500038}{doi:10.1142/S1230161223500038};
\href{https://arxiv.org/abs/2211.08351v3}{arXiv:2211.08351v3}.
\bibitem{ss} M. Skeide and K. Sumesh, CP-H-extendable maps between
Hilbert modules and CPH-semigroups, J. Math. Anal. Appl. \textbf{414}
(2014), 886--913.
\href{https://doi.org/10.1016/j.jmaa.2014.01.024}{doi:10.1016/j.jmaa.2014.01.024};
\href{https://arxiv.org/abs/1210.7491v2}{arXiv:1210.7491v2}.
\bibitem{shalit} O. M. Shalit and M. Skeide, CP-Semigroups and
Dilations, Subproduct Systems and Superproduct Systems: The Multi-Parameter
Case and Beyond, Dissertationes Math. \textbf{585} (2023).
\href{https://doi.org/10.4064/dm823-5-2022}{doi:10.4064/dm823-5-2022};
\href{https://arxiv.org/abs/2003.05166v3}{arXiv:2003.05166v3}.
\end{thebibliography}
\end{document}
